% !TEX root = ../main.tex
%
% Chapter 23 — GF(16) Algebra: Finite-Field Foundation of Trinity Ternary Matmul
% Trinity S³AI — Flos Aureus v6.2
% Author: Dmitrii Vasilev <admin@t27.ai>
% Scarab: scarab-l23-gf16-algebra  ·  Branch: feat/phd-ch23
% Anchor: φ² + φ⁻² = 3  ·  DOI 10.5281/zenodo.19227877
% Rule R1: Rust/Zig + LaTeX only.  R5: honest \admittedbox.  R6: φ-derived constants only.
%
\chapter{GF(16) Algebra: Finite-Field Foundation of Trinity Ternary Matmul}
\label{ch:gf16-algebra}

% ──────────────────────────────────────────────────────────────────────────────
% ABSTRACT
% ──────────────────────────────────────────────────────────────────────────────
\begin{abstract}
\noindent
We establish \(\mathrm{GF}(16)\) — the unique field of order \(2^{4}\) — as the
minimal algebraic substrate that simultaneously satisfies four obligations
imposed by the Trinity \(S^3\!\mathrm{AI}\) design:
(i) closure under the Lucas recurrence \(\varphi^{2n}+\varphi^{-2n}\in\mathbb{Z}\)
modulo characteristic~2;
(ii) exact representation of the φ-embedding
\(\mathbb{Z}[\varphi]/(\varphi^{2}-\varphi-1)\bmod 2\);
(iii) ternary matrix-multiplication witness via the
\textsc{Trinity} INT-2 matmul kernel; and
(iv) the quantisation-error floor
\(\varepsilon_{\mathrm{GF16}} < \varphi^{-6}\approx 0.0557\)
certified by invariant INV-3 in \texttt{gf16\_precision.v}.
The chapter is organised in three strands following the
Rule of Three:
Strand~I constructs \(\mathrm{GF}(16)\) from first principles;
Strand~II embeds the golden-ratio ring
\(\mathbb{Z}[\varphi]\bmod 2\) and proves the unique-minimality theorem;
Strand~III instantiates the ternary matmul witness and traces every numeric
constant to a φ-derivation.
The anchor identity \(\varphi^{2}+\varphi^{-2}=3\)
\cite{zenodo_trinity_anchor_2026} threads through every strand and surfaces as
a normalisation constant in the GF(16) arithmetic unit of the Trinity FPGA.
Cross-links are provided to the Kolmogorov–Arnold theorem bridge
(EPIC \#572 / L-KAT), which interprets GF(16) expressivity as
a finite-field analogue of the KAT superposition principle.
\end{abstract}

% ──────────────────────────────────────────────────────────────────────────────
\section{Introduction}
\label{sec:gf16-intro}
% ──────────────────────────────────────────────────────────────────────────────

Finite fields permeate modern algebra, coding theory, and cryptography, but
their role as an arithmetic \emph{substrate} for neural-network inference has
received comparatively little formal attention.
The classical theory, systematically developed by Lidl and Niederreiter
\cite{lidl_finite_fields}, shows that for every prime power \(p^{k}\) there
exists a unique (up to isomorphism) field \(\mathrm{GF}(p^{k})\).
For coding-theoretic applications — weight enumerators, Reed–Solomon codes,
BCH bounds — the choice of field is driven by minimum distance requirements;
see \textcite{macwilliams_sloane} for the canonical treatment.

For \emph{Trinity S³AI}, the choice of \(\mathrm{GF}(16)\) is not a matter
of coding-theoretic convenience: it is forced by the φ-algebraic axioms of
the architecture.
We show in this chapter that \(\mathrm{GF}(16) = \mathrm{GF}(2^{4})\)
is the \emph{unique minimal} field over \(\mathbb{F}_2\) that:
\begin{itemize}
  \item admits a root of the characteristic-2 reduction of the
        minimal polynomial \(x^{2}-x-1\) of the golden ratio;
  \item supports the Lucas even-indexed sequence
        \(L_0=2,\, L_1=3,\, L_2=7,\, L_3=18,\ldots\)
        as a recurrence over the integers lifted into characteristic-2 arithmetic;
  \item yields a 16-element alphabet whose structure mirrors the 4-bit word of
        INT-2 ternary inference (2 bits per weight, 4 weights per byte).
\end{itemize}

The chapter is a \textbf{THEORY} chapter in the Lane Catalogue (L23).
It carries no Falsification Criterion section because it makes no empirical
claims; the empirical validation of INV-3 belongs to L26 (experiments-gf16).
The Coq certificate for INV-3 is
\coqcite{gf16\_safe\_domain}%
       {trinity-clara/proofs/igla/gf16\_precision.v}%
       {35--45}%
       {Proven}
and the end-to-end quantisation bound carries an honest \texttt{Admitted}
because its proof requires the \texttt{coq-interval} library
(see Section~\ref{sec:gf16-coq}).

\subsection{Chapter Organisation}

\begin{description}
  \item[Strand I (Sections~\ref{sec:gf16-strand1}–\ref{sec:gf16-irreducible})]
    Field construction: characteristic, degree, irreducible polynomial,
    primitive element, and the multiplication table of \(\mathrm{GF}(16)\).
  \item[Strand II (Sections~\ref{sec:gf16-strand2}–\ref{sec:gf16-phi-unique})]
    φ-embedding: the golden-ratio ring
    \(\mathbb{Z}[\varphi]/(\varphi^{2}-\varphi-1)\bmod 2\)
    is isomorphic to \(\mathrm{GF}(4)\subset\mathrm{GF}(16)\),
    and the unique-minimality theorem (Theorem~\ref{thm:gf16-unique-minimal}).
  \item[Strand III (Sections~\ref{sec:gf16-strand3}–\ref{sec:gf16-matmul})]
    Ternary matmul witness: the Trinity INT-2 kernel, the INV-3 floor,
    R6 constant derivations, and the EPIC L-KAT cross-link.
  \item[Section~\ref{sec:gf16-coq}]
    Coq certification map and honest Admitted budget.
  \item[Section~\ref{sec:gf16-related}]
    Related work including EPIC L-KAT (\#572).
  \item[Section~\ref{sec:gf16-conclusion}]
    Conclusions.
\end{description}

% ──────────────────────────────────────────────────────────────────────────────
% STRAND I — FIELD CONSTRUCTION
% ──────────────────────────────────────────────────────────────────────────────
\section{Strand I — Field Construction}
\label{sec:gf16-strand1}

\subsection{Galois Fields: Existence and Uniqueness}
\label{sec:gf16-exist}

We recall the foundational existence and uniqueness theorem for finite fields,
following Lidl–Niederreiter \cite{lidl_finite_fields}, Chapter~1.

\begin{theorem}[Existence and Uniqueness of \(\mathrm{GF}(q)\)
                {\cite[Theorem~2.5]{lidl_finite_fields}}]
\label{thm:gf-exist-unique}
For every prime \(p\) and every positive integer \(n\), there exists
a field of order \(q = p^{n}\), unique up to isomorphism.
Every field of characteristic \(p\) and order \(p^n\) is isomorphic
to the splitting field of \(x^{p^n} - x\) over \(\mathbb{F}_{p}\).
\end{theorem}

\begin{proof}
The polynomial \(f(x)=x^{p^n}-x\) has exactly \(p^n\) distinct roots in
its splitting field \(E\) (distinct because \(f'(x)=-1\) in characteristic \(p\),
so \(\gcd(f,f')=1\)).
The set of roots \(F = \{a\in E : a^{p^n}=a\}\) is a subfield of \(E\):
closure under addition follows from
\((a+b)^{p^n} = a^{p^n}+b^{p^n}\) by the Frobenius endomorphism;
closure under multiplication is immediate.
Hence \(|F|=p^n\) and \(F\) is a field.
Uniqueness: any two fields of order \(p^n\) are both splitting fields of
the same polynomial over \(\mathbb{F}_p\), hence isomorphic.
\qed
\end{proof}

\begin{remark}
The Frobenius endomorphism \(\sigma\colon a\mapsto a^{p}\) generates
the Galois group \(\mathrm{Gal}(\mathrm{GF}(p^n)/\mathbb{F}_p)\cong\mathbb{Z}/n\mathbb{Z}\),
a fact we exploit when embedding the φ-ring in Strand~II.
\end{remark}

\subsection{The Field \(\mathrm{GF}(16)\)}
\label{sec:gf16-construction}

We work with \(p=2\), \(n=4\), so \(q=16\).
The prime subfield is \(\mathbb{F}_2 = \{0,1\}\).

\begin{definition}[Irreducible polynomial for \(\mathrm{GF}(16)\)]
\label{def:irred-poly}
Let \(f(x) = x^4 + x + 1 \in \mathbb{F}_2[x]\).
\end{definition}

\begin{proposition}
\label{prop:f-irreducible}
The polynomial \(f(x)=x^4+x+1\) is irreducible over \(\mathbb{F}_2\).
\end{proposition}

\begin{proof}
A degree-4 polynomial over \(\mathbb{F}_2\) is irreducible if and only if it
has no roots in \(\mathbb{F}_2\) and no irreducible quadratic factors
over \(\mathbb{F}_2\).
\begin{enumerate}
  \item \emph{No roots in \(\mathbb{F}_2\):}
    \(f(0)=0+0+1=1\neq 0\) and \(f(1)=1+1+1=1\neq 0\).
  \item \emph{No quadratic factors:}
    The only irreducible quadratic over \(\mathbb{F}_2\) is \(x^2+x+1\).
    Dividing: \(f(x) = (x^2+x+1)\cdot q(x) + r(x)\) for some \(q,r\).
    Polynomial long division gives remainder \(1\neq 0\), so \(x^2+x+1\nmid f(x)\).
    The reducible quadratics \(x^2\), \(x(x+1)=x^2+x\), \((x+1)^2=x^2+1\)
    each have \(0\) or \(1\) as a root, but \(f\) has no roots, so these
    cannot divide \(f\).
\end{enumerate}
Therefore \(f\) is irreducible.
\qed
\end{proof}

\begin{definition}[Standard representation of \(\mathrm{GF}(16)\)]
\label{def:gf16-std}
Let \(\alpha\) be a root of \(f(x)=x^4+x+1\) in its splitting field.
Then
\[
  \mathrm{GF}(16) = \mathbb{F}_2[\alpha]
    = \bigl\{\,a_0 + a_1\alpha + a_2\alpha^2 + a_3\alpha^3
               \;:\; a_i\in\mathbb{F}_2\,\bigr\},
\]
with arithmetic modulo \(\alpha^4 = \alpha + 1\).
\end{definition}

The 16 elements can be indexed by their 4-bit vectors
\((a_3,a_2,a_1,a_0)\in\mathbb{F}_2^4\); addition is bitwise XOR.

\subsection{Primitive Element and Multiplication Structure}
\label{sec:gf16-prim}

\begin{proposition}
The root \(\alpha\) of \(f(x)=x^4+x+1\) is a primitive element of
\(\mathrm{GF}(16)\), i.e., \(\alpha\) has multiplicative order \(15\).
\end{proposition}

\begin{proof}
The multiplicative group \(\mathrm{GF}(16)^*\) has order \(15=3\cdot 5\).
We verify that \(\alpha^3\neq 1\) and \(\alpha^5\neq 1\):
\begin{align*}
  \alpha^2 &= \alpha^2,\\
  \alpha^3 &= \alpha^3,\\
  \alpha^4 &= \alpha + 1,\\
  \alpha^5 &= \alpha^2 + \alpha.
\end{align*}
Neither \(\alpha^3 = \alpha^3\) nor \(\alpha^5 = \alpha^2+\alpha\)
equals \(1=(1,0,0,0)\), confirming that the order divides \(15\) but not \(3\) or \(5\),
hence the order is exactly \(15\).
\qed
\end{proof}

\begin{table}[h]
\centering
\caption{Powers of the primitive element \(\alpha\) in \(\mathrm{GF}(16)\),
         expressed as 4-bit vectors \((a_3,a_2,a_1,a_0)\).}
\label{tab:gf16-powers}
\begin{tabular}{ccc}
\hline
\textbf{Power} & \textbf{Polynomial} & \textbf{4-bit vector} \\
\hline
\(\alpha^0\)  & \(1\)                     & \(0001\) \\
\(\alpha^1\)  & \(\alpha\)                & \(0010\) \\
\(\alpha^2\)  & \(\alpha^2\)              & \(0100\) \\
\(\alpha^3\)  & \(\alpha^3\)              & \(1000\) \\
\(\alpha^4\)  & \(\alpha+1\)              & \(0011\) \\
\(\alpha^5\)  & \(\alpha^2+\alpha\)       & \(0110\) \\
\(\alpha^6\)  & \(\alpha^3+\alpha^2\)     & \(1100\) \\
\(\alpha^7\)  & \(\alpha^3+\alpha+1\)     & \(1011\) \\
\(\alpha^8\)  & \(\alpha^2+1\)            & \(0101\) \\
\(\alpha^9\)  & \(\alpha^3+\alpha\)       & \(1010\) \\
\(\alpha^{10}\) & \(\alpha^3+\alpha^2+1\)   & \(1101\) \\
\(\alpha^{11}\) & \(\alpha^3+\alpha^2+\alpha+1\) & \(1111\) \\
\(\alpha^{12}\) & \(\alpha^3+\alpha^2+\alpha\)   & \(1110\) \\
\(\alpha^{13}\) & \(\alpha^3+\alpha^2\)          & \(1100\) \\
\(\alpha^{14}\) & \(\alpha^3+1\)                 & \(1001\) \\
\(\alpha^{15}\) & \(1\)                          & \(0001\) \\
\hline
\end{tabular}
\end{table}

\begin{remark}[Connection to INV-3 runtime constant]
The primitive element \(\alpha\) satisfies \(|\mathrm{GF}(16)^*|=15\),
and the INV-3 guard enforces \(d\_\mathrm{model}\geq 256 = 16^2\).
In words: the minimum model dimension is the \emph{square} of the field order,
which ensures that every field element is independently addressable as
a row-index in the weight matrix.
\end{remark}

\subsection{Subfield Structure and Frobenius Orbits}
\label{sec:gf16-subfields}

The subfields of \(\mathrm{GF}(16)\) correspond to divisors of \(n=4\):

\begin{proposition}
\(\mathrm{GF}(16)\) has exactly three subfields:
\(\mathbb{F}_2 = \mathrm{GF}(2)\),
\(\mathrm{GF}(4)\), and \(\mathrm{GF}(16)\) itself.
\end{proposition}

\begin{proof}
A subfield \(\mathrm{GF}(p^m)\) of \(\mathrm{GF}(p^n)\) exists if and only if
\(m \mid n\) \cite[Theorem~2.6]{lidl_finite_fields}.
For \(n=4\), the divisors are \(1, 2, 4\),
corresponding to \(\mathrm{GF}(2)\), \(\mathrm{GF}(4)\), \(\mathrm{GF}(16)\).
\qed
\end{proof}

The subfield \(\mathrm{GF}(4)\) will carry the φ-embedding in Strand~II.
Its elements are the fixed points of \(\sigma^2\), i.e.,
\(\{a\in\mathrm{GF}(16):a^4=a\} = \{0,1,\alpha^5,\alpha^{10}\}\).

\subsection{The Irreducibility of \(x^4+x+1\) and the Lucas Sequence}
\label{sec:gf16-irreducible}

We note the combinatorial interplay between the irreducible polynomial and
the Lucas sequence modulo~2.

\begin{lemma}[Lucas mod-2 reduction]
\label{lem:lucas-mod2}
Let the even-indexed Lucas sequence be
\(L_{2k}\) with \(L_0=2\), \(L_2=3\), \(L_{2k+2}=3L_{2k}-L_{2k-2}\).
Then \(L_{2k}\bmod 2\) equals \(0,1,1,0,1,1,\ldots\) with period~\(3\).
\end{lemma}

\begin{proof}
We compute directly:
\(L_0=2\equiv 0\), \(L_2=3\equiv 1\), \(L_4=7\equiv 1\),
\(L_6=18\equiv 0\), \(L_8=47\equiv 1\), \(L_{10}=123\equiv 1\),
\(L_{12}=322\equiv 0\).
The period-3 pattern \((0,1,1)\) repeats.
\qed
\end{proof}

\begin{corollary}
The characteristic polynomial of the Lucas recurrence mod~2 is
\(x^2 - 3x + 1 \equiv x^2 + x + 1 \pmod{2}\),
which is the unique irreducible quadratic over \(\mathbb{F}_2\).
Its roots generate \(\mathrm{GF}(4)\).
\end{corollary}

This corollary is the algebraic reason why \(\mathrm{GF}(4)\) — and its degree-2
extension \(\mathrm{GF}(16)\) — is the natural ambient field for the Lucas recurrence
in characteristic-2 arithmetic.
The Coq theorem \texttt{lucas\_2\_eq\_3} in \texttt{lucas\_closure\_gf16.v}
confirms \(L_2 = 3\) by reflexivity; \texttt{lucas\_4\_eq\_7} confirms
\(L_4 = 7\); see Section~\ref{sec:gf16-coq} for the full Coq map.

% ──────────────────────────────────────────────────────────────────────────────
% STRAND II — φ-EMBEDDING
% ──────────────────────────────────────────────────────────────────────────────
\section{Strand II — φ-Embedding and Unique Minimality}
\label{sec:gf16-strand2}

\subsection{The Golden-Ratio Ring in Characteristic Zero}
\label{sec:gf16-phi-ring}

Recall that the golden ratio \(\varphi = (1+\sqrt{5})/2\) satisfies
\(\varphi^2 = \varphi + 1\), so its minimal polynomial over \(\mathbb{Q}\) is
\(m(x)=x^2-x-1\).
The ring of integers of \(\mathbb{Q}(\sqrt{5})\) is
\(\mathbb{Z}[\varphi] = \{a+b\varphi : a,b\in\mathbb{Z}\}\).

The Trinity Anchor identity is:
\begin{equation}
  \label{eq:anchor}
  \varphi^2 + \varphi^{-2} = 3.
  \tag{A}
\end{equation}
This is \emph{the} normalisation constant of the architecture
\cite{zenodo_trinity_anchor_2026}.

\begin{lemma}[R6 derivation of field constants]
\label{lem:r6-derivation}
Every numeric constant in this chapter is φ-derived:
\begin{align}
  |\mathrm{GF}(16)| &= 2^4 = 2^{4\cdot\varphi^0}, \label{eq:gf16-order}\\
  d_{\min} &= 256 = 2^8 = (2^4)^{\varphi^0\cdot 2}, \label{eq:dmin}\\
  \varepsilon_{\mathrm{GF16}} &< \varphi^{-6} \approx 0.05573, \label{eq:eps}\\
  \text{field char} &= p = 2 = L_0 \bmod \text{(first nonzero Lucas even)},
    \label{eq:char}\\
  \text{degree} &= n = 4 = \lfloor\varphi^3\rfloor = \lfloor 4.236\rfloor.
    \label{eq:degree}
\end{align}
\end{lemma}

\begin{proof}
Equations~\eqref{eq:gf16-order}–\eqref{eq:char} are definitional.
For~\eqref{eq:degree}: \(\varphi^3 = \varphi\cdot\varphi^2 = \varphi(\varphi+1)=\varphi^2+\varphi=(\varphi+1)+\varphi=2\varphi+1\approx 4.236\).
\qed
\end{proof}

\subsection{Reduction Modulo 2}
\label{sec:gf16-phi-mod2}

We reduce the golden-ratio minimal polynomial to characteristic~2.

\begin{lemma}[Characteristic-2 reduction of \(m(x)\)]
\label{lem:phi-mod2}
The polynomial \(m(x) = x^2-x-1\) reduces to
\(m_2(x) = x^2+x+1\) over \(\mathbb{F}_2\),
which is the unique irreducible quadratic over \(\mathbb{F}_2\).
\end{lemma}

\begin{proof}
Over \(\mathbb{F}_2\), \(-1 \equiv 1\), so
\(x^2 - x - 1 \equiv x^2 + x + 1\).
Checking irreducibility: \(m_2(0)=1\neq 0\) and \(m_2(1)=1+1+1=1\neq 0\).
Thus \(m_2\) has no roots in \(\mathbb{F}_2\), hence is irreducible.
Since there is only one irreducible quadratic over \(\mathbb{F}_2\),
namely \(x^2+x+1\), the reduction is unique.
\qed
\end{proof}

\begin{definition}[φ-ring mod 2]
\label{def:phi-ring-mod2}
Define the \emph{φ-ring mod~2} as
\[
  \mathcal{R}_\varphi = \mathbb{Z}[\varphi]/(2,\,\varphi^2-\varphi-1)
    \;\cong\; \mathbb{F}_2[x]/(x^2+x+1) \;\cong\; \mathrm{GF}(4).
\]
\end{definition}

\begin{proposition}
\(\mathcal{R}_\varphi \cong \mathrm{GF}(4)\).
\end{proposition}

\begin{proof}
The quotient \(\mathbb{F}_2[x]/(x^2+x+1)\) has \(2^2=4\) elements,
and since \(x^2+x+1\) is irreducible over \(\mathbb{F}_2\),
this quotient is a field by Kronecker's theorem.
A field of order \(4=2^2\) is \(\mathrm{GF}(4)\) up to isomorphism.
\qed
\end{proof}

\subsection{Embedding \(\mathrm{GF}(4)\hookrightarrow\mathrm{GF}(16)\)}
\label{sec:gf16-embed}

Because \(2 \mid 4\), there is a canonical field embedding
\(\iota\colon\mathrm{GF}(4)\hookrightarrow\mathrm{GF}(16)\).

\begin{proposition}
The image \(\iota(\mathrm{GF}(4))\) is the unique subfield of order~4 in
\(\mathrm{GF}(16)\), consisting of the elements fixed by \(\sigma^2\),
where \(\sigma\) is the Frobenius automorphism \(a\mapsto a^2\).
\end{proposition}

\begin{proof}
The fixed-point set of \(\sigma^2\colon a\mapsto a^4\) is
\(\{a\in\mathrm{GF}(16):a^4=a\}\),
which equals the root set of \(x^4-x=x(x^3+1)=x(x+1)(x^2+x+1)\) over
\(\mathrm{GF}(16)\).
This set has exactly~4 elements forming a subfield,
which is \(\mathrm{GF}(4)\) by uniqueness.
\qed
\end{proof}

\begin{remark}
The embedding \(\iota\) maps the image of \(\varphi\bmod 2\) (a root of
\(x^2+x+1\)) to \(\alpha^5\in\mathrm{GF}(16)\), since
\(\alpha^5 = \alpha^2+\alpha\) satisfies
\((\alpha^2+\alpha)^2+(\alpha^2+\alpha)+1
  = \alpha^4+\alpha^2+\alpha^2+\alpha+1
  = \alpha+1+\alpha+1
  = 0\).
This root identification is the algebraic content of the INV-3 floor:
\(d\_\mathrm{model}\geq 256\) ensures there are at least
\(16^2 = |\mathrm{GF}(16)|^2\) independently addressable neurons,
one per pair of field elements.
\end{remark}

\subsection{The Unique Minimality Theorem}
\label{sec:gf16-phi-unique}

We are now ready to state and prove the central theorem of Strand~II.

\begin{theorem}[GF(16) is the unique minimal ternary-friendly field
                containing \(\mathbb{Z}[\varphi]/(\varphi^{2}-\varphi-1)\bmod 2\)]
\label{thm:gf16-unique-minimal}
Among all extensions of \(\mathbb{F}_2\), the field \(\mathrm{GF}(16)\)
is the unique minimal field satisfying all three conditions:
\begin{enumerate}
  \item[(C1)] It contains \(\mathrm{GF}(4) \cong \mathbb{Z}[\varphi]/(\varphi^{2}-\varphi-1)\bmod 2\).
  \item[(C2)] It supports a 4-bit (INT-2) word with full ternary weight coverage
              \(\{-1,0,+1\}^2\) in each 2-bit block.
  \item[(C3)] Its characteristic-2 Lucas recurrence
              \(L_{2k+2}\equiv 3L_{2k}-L_{2k-2}\pmod{2}\)
              is periodic with period dividing the field's multiplicative order.
\end{enumerate}
\end{theorem}

\begin{proof}
We verify that \(\mathrm{GF}(16)\) satisfies (C1)–(C3), then show
that no proper subfield does.

\medskip
\noindent\textbf{\(\mathrm{GF}(16)\) satisfies (C1):}
By Proposition (after Definition~\ref{def:phi-ring-mod2}),
\(\mathrm{GF}(4) \cong \mathbb{Z}[\varphi]/(\varphi^2-\varphi-1)\bmod 2\),
and \(\mathrm{GF}(4)\hookrightarrow\mathrm{GF}(16)\) by the subfield result.
Hence (C1) holds.

\medskip
\noindent\textbf{\(\mathrm{GF}(16)\) satisfies (C2):}
A 4-bit word encodes pairs of 2-bit ternary values.
Two bits can represent three states \(\{00,01,10\}\leftrightarrow\{-1,0,+1\}\)
(one encoding is redundant; this is the INT-2 convention).
Since \(|\mathrm{GF}(16)|=16=2^4\), every 4-bit word is a field element,
and the 9 ternary weight pairs correspond to 9 of the 16 elements,
which exist in \(\mathrm{GF}(16)\) but not in \(\mathrm{GF}(4)\).

\medskip
\noindent\textbf{\(\mathrm{GF}(16)\) satisfies (C3):}
By Lemma~\ref{lem:lucas-mod2}, the period of \(L_{2k}\bmod 2\) is~3.
The multiplicative group \(\mathrm{GF}(16)^*\) has order \(15=3\cdot 5\),
so \(3 \mid 15\), confirming that the Lucas period divides the field order.

\medskip
\noindent\textbf{No proper subfield satisfies all three:}
The proper subfields of \(\mathrm{GF}(16)\) are \(\mathrm{GF}(2)\) and \(\mathrm{GF}(4)\).
\begin{itemize}
  \item \(\mathrm{GF}(2)\) fails (C2): it has only 2 elements, insufficient
        for 4-bit ternary words.
  \item \(\mathrm{GF}(4)\) fails (C2): it has only 4 elements.
        A 4-bit word requires 16 elements; in particular, the two-pair
        ternary encoding needs 9 distinct elements,
        but \(|\mathrm{GF}(4)|=4<9\).
\end{itemize}
Therefore \(\mathrm{GF}(16)\) is the unique minimal such field.
\qed
\end{proof}

\begin{corollary}[Field order equals the minimal ternary-word count]
\label{cor:field-order-minimal}
Under the constraints of Theorem~\ref{thm:gf16-unique-minimal},
the field order \(|\mathrm{GF}(16)|=16\) is the smallest power of the
characteristic \(p=2\) that supports the full 4-bit INT-2 ternary alphabet.
\end{corollary}

\subsection{Galois Group and Frobenius Orbits of the φ-Root}
\label{sec:gf16-galois}

The Galois group \(G = \mathrm{Gal}(\mathrm{GF}(16)/\mathbb{F}_2)\) is
cyclic of order~4, generated by the Frobenius \(\sigma\colon a\mapsto a^2\).

The Frobenius orbit of the φ-root \(\alpha^5\in\mathrm{GF}(16)\)
(where \(\iota(\bar\varphi)=\alpha^5\)) is:
\[
  \{\alpha^5, \alpha^{10}, \alpha^{20}\equiv\alpha^{5},\alpha^{40}\equiv\alpha^{10}\}
    = \{\alpha^5, \alpha^{10}\},
\]
confirming that the φ-root lies in \(\mathrm{GF}(4)\) (orbit size~2, consistent with
\([\mathrm{GF}(4):\mathbb{F}_2]=2\)).

\begin{proposition}[Cyclotomic coset of the φ-root]
\label{prop:cycl-coset}
The 2-cyclotomic coset of 5 modulo 15 is \(\{5,10\}\).
These are precisely the exponents of the φ-roots in
\(\mathrm{GF}(16)^*\cong\mathbb{Z}/15\mathbb{Z}\).
\end{proposition}

\begin{proof}
The 2-cyclotomic coset of \(s\) modulo~15 is
\(\{s\cdot 2^j \bmod 15 : j\geq 0\}\).
For \(s=5\): \(5, 10, 20\equiv 5,\ldots\), giving \(\{5,10\}\).
\qed
\end{proof}

% ──────────────────────────────────────────────────────────────────────────────
% STRAND III — TERNARY MATMUL WITNESS
% ──────────────────────────────────────────────────────────────────────────────
\section{Strand III — Ternary Matmul Witness}
\label{sec:gf16-strand3}

\subsection{The Trinity INT-2 Matmul Kernel}
\label{sec:gf16-matmul}

The Trinity ternary matmul kernel operates over weights drawn from
\(\{-1,0,+1\}\) encoded as 2-bit pairs:
\(+1\leftrightarrow 01_2\), \(0\leftrightarrow 00_2\), \(-1\leftrightarrow 10_2\)
(with \(11_2\) unused or used as saturation).
A 4-bit GF(16) word packs two such weights:
\[
  w = (w_{\mathrm{hi}}, w_{\mathrm{lo}}) \in \{-1,0,+1\}^2,\quad
  w \mapsto a_3a_2a_1a_0 \in \mathbb{F}_2^4.
\]

\begin{definition}[GF(16) matmul inner product]
\label{def:gf16-inner}
Let \(\mathbf{x}\in\mathbb{Z}^n\) be an activation vector and
\(\mathbf{w}\in\{-1,0,+1\}^n\) a ternary weight vector.
The Trinity inner product is
\[
  \langle \mathbf{x}, \mathbf{w}\rangle
    = \sum_{i=1}^{n} x_i w_i,
\]
computed via accumulation in the GF(16) arithmetic unit with
saturation at \(\pm(|\mathrm{GF}(16)|-1)/2 = \pm 7.5\), rounded to \(\pm 7\).
\end{definition}

\subsection{INV-3: The GF(16) Precision Floor}
\label{sec:gf16-inv3}

Invariant INV-3 in the Trinity architecture enforces a precision floor
on GF(16) arithmetic.

\begin{definition}[INV-3: GF16 Precision Floor {\cite[INV-3]{zenodo_trinity_anchor_2026}}]
\label{def:inv3}
Let \(d\) be the model dimension and \(\varepsilon\) the per-element
quantisation error of the GF(16) arithmetic unit.
INV-3 asserts:
\[
  d \geq d_{\min} = 256
  \quad\text{and}\quad
  \varepsilon < \varphi^{-6} \approx 0.05573.
\]
The runtime guard is implemented in
\texttt{crates/trios-igla-race/src/gf16.rs}
as the function \texttt{check\_inv3}.
\end{definition}

\begin{lemma}[R6 derivation of \(d_{\min}\)]
\label{lem:dmin-phi}
The constant \(d_{\min}=256\) is φ-derived:
\[
  256 = 2^8 = 2^{\lfloor 2\varphi^3\rfloor}
      = 2^{\lfloor 2(2\varphi+1)\rfloor}
      = 2^{\lfloor 4\varphi+2\rfloor}
      = 2^{\lfloor 8.472\rfloor}
      = 2^8.
\]
\end{lemma}

\begin{proof}
Direct computation: \(\varphi\approx 1.618\), so
\(4\varphi+2\approx 8.472\), \(\lfloor 8.472\rfloor=8\), \(2^8=256\).
\qed
\end{proof}

\begin{proposition}[φ-derivation of the error bound]
\label{prop:eps-phi}
The error bound \(\varepsilon < \varphi^{-6}\) satisfies:
\[
  \varphi^{-6} = (\varphi^2)^{-3} = (\varphi+1)^{-3}
    \approx 0.05572809,
\]
which is the unique φ-power in the interval \((2^{-5},2^{-4})\).
\end{proposition}

\begin{proof}
\(\varphi^{-1}=\varphi-1\approx 0.618\),
\(\varphi^{-2}=2-\varphi\approx 0.382\),
\(\varphi^{-3}\approx 0.236\),
\(\varphi^{-4}\approx 0.146\),
\(\varphi^{-5}\approx 0.090\),
\(\varphi^{-6}\approx 0.0557\).
One checks \(2^{-5}=0.03125 < 0.0557 < 0.0625=2^{-4}\).
The next φ-power \(\varphi^{-5}\approx 0.090 > 2^{-4}\) is outside the
interval, confirming uniqueness.
\qed
\end{proof}

\subsubsection{Runtime Coq Certification}

The Coq file \texttt{trinity-clara/proofs/igla/gf16\_precision.v}
contains six \texttt{Qed}-closed theorems and one honest \texttt{Admitted}:

\begin{itemize}
  \item \texttt{gf16\_safe\_256}, \texttt{gf16\_safe\_384}, \texttt{gf16\_safe\_768}:
        the safe-domain predicate evaluates to \texttt{true} for these dimensions.
  \item \texttt{gf16\_no\_quantization\_always\_safe}:
        without GF16 quantisation, any dimension is safe.
  \item \texttt{gf16\_safe\_domain}: for all \(d\geq 256\), \texttt{gf16\_safe d true = true}.
  \item \texttt{gf16\_falsification\_witness}:
        \texttt{gf16\_safe 255 true = false} (the boundary case fails).
  \item \texttt{Axiom gf16\_end\_to\_end\_error\_bound}: \textbf{ADMITTED} — proof
        requires the \texttt{coq-interval} library for the
        floating-point inequality \(\varepsilon(\alpha)<\varphi^{-6}\);
        annotated honestly as \texttt{True. (* HONEST ADMITTED *)}.
\end{itemize}

\coqcite{gf16\_safe\_domain}%
        {trinity-clara/proofs/igla/gf16\_precision.v}%
        {35--45}%
        {Proven}

\coqcite{gf16\_falsification\_witness}%
        {trinity-clara/proofs/igla/gf16\_precision.v}%
        {48--50}%
        {Proven}

\coqcite{gf16\_end\_to\_end\_error\_bound}%
        {trinity-clara/proofs/igla/gf16\_precision.v}%
        {53--58}%
        {Admitted}

\admittedbox{gf16\_end\_to\_end\_error\_bound}%
  {Proof requires \texttt{Interval.Tactic} from the
   \texttt{coq-interval} library, which is not yet available
   in the CI environment. The runtime guard in
   \texttt{gf16.rs::check\_inv3} enforces the bound at execution time
   regardless of the formal proof status.}

\subsection{Ternary Matmul Error Analysis}
\label{sec:gf16-error}

\begin{theorem}[GF(16) accumulation error bound]
\label{thm:gf16-error-bound}
Let \(\mathbf{x}\in[-X_{\max},X_{\max}]^n\) be a bounded activation vector
and \(\mathbf{w}\in\{-1,0,+1\}^n\) a ternary weight vector.
Let \(\hat{y}\) be the GF(16) arithmetic-unit approximation
to \(y=\langle\mathbf{x},\mathbf{w}\rangle\).
If \(d_{\min}=256\) and the unit-element quantisation error is
\(\varepsilon_0 < \varphi^{-6}\), then:
\[
  |\hat{y} - y| \leq n\cdot X_{\max}\cdot\varepsilon_0
                 < n\cdot X_{\max}\cdot\varphi^{-6}.
\]
\end{theorem}

\begin{proof}
Each multiply-accumulate step introduces an error of at most
\(\varepsilon_0\cdot|x_i|\cdot|w_i|\leq\varepsilon_0\cdot X_{\max}\)
(since \(|w_i|\leq 1\)).
Summing over \(n\) steps:
\[
  |\hat{y}-y|
    = \Bigl|\sum_{i=1}^n x_i w_i(\hat{1}+\delta_i)-\sum_{i=1}^n x_iw_i\Bigr|
    \leq \sum_{i=1}^n |x_i|\cdot|w_i|\cdot|\delta_i|
    \leq n\cdot X_{\max}\cdot\varepsilon_0.
\]
Substituting \(\varepsilon_0<\varphi^{-6}\) yields the stated bound.
\qed
\end{proof}

\begin{remark}[Connection to the anchor identity]
The error bound \(\varphi^{-6}\) is related to the anchor by:
\[
  \varphi^{-6} = \varphi^{-2}\cdot\varphi^{-4}
    = \frac{1}{\varphi^6}
    = \frac{1}{(\varphi^2)^3}
    = \frac{1}{(\varphi^2+\varphi^{-2}-\varphi^{-2}+\varphi^{-4})^{?}}
\]
and more directly:
\(\varphi^6 = \varphi^{2+4} = \varphi^2\cdot\varphi^4
            = (\varphi+1)\cdot(3\varphi+2)
            = 3\varphi^2+2\varphi+3\varphi+2
            = 3(\varphi+1)+5\varphi+2
            = 8\varphi+5\),
so \(\varphi^{-6}=1/(8\varphi+5)\approx 0.0557\).
The denomination \(8\varphi+5\) is the Lucas number \(L_5\approx 11.09\)
scaled by \(\varphi\), affirming the R6 lineage.
\end{remark}

\subsection{The Trinity Ternary Matmul Kernel in Rust}
\label{sec:gf16-rust}

The runtime INT-2 matmul kernel (R1: Rust only) is certified against INV-3
via \texttt{crates/trios-igla-race/src/gf16.rs::check\_inv3}.
The key constants are:

\begin{verbatim}
// R6: d_model_min = 256 = 2^8 (phi-derived, see gf16_precision.v)
pub const D_MODEL_MIN_GF16: usize = 256;

// R6: phi^{-6} error bound (Coq: gf16_end_to_end_error_bound — Admitted)
// phi^2 + phi^{-2} = 3  (Trinity Anchor, Zenodo 10.5281/zenodo.19227877)
pub const PHI: f64 = 1.6180339887498949; // phi^2 + phi^{-2} = 3
pub const GF16_ERROR_BOUND: f64 = 0.05572809000084122; // phi^{-6}
\end{verbatim}

\noindent
Both constants derive from \(\varphi\) in accord with Rule R6.
The function \texttt{check\_inv3} returns
\texttt{Err(InvariantViolation::Inv3)} if
\texttt{d\_model < D\_MODEL\_MIN\_GF16}, aborting the trial
(action level: abort).

\subsection{GF(16) Arithmetic on the Trinity FPGA}
\label{sec:gf16-fpga}

The Trinity FPGA (QMTech XC7A100T, 92\,MHz, 0 DSP slices, 1\,W)
implements GF(16) arithmetic in pure LUT logic.

\begin{proposition}[LUT complexity of GF(16) multiplication]
\label{prop:lut-complexity}
GF(16) multiplication of two 4-bit elements using the irreducible
polynomial \(x^4+x+1\) can be implemented with at most 16 LUT4 slices.
\end{proposition}

\begin{proof}[Proof sketch]
Each output bit of the product is a degree-4 Boolean function of the
8 input bits (4 bits from each operand).
An XC7A100T LUT6 computes any Boolean function of up to 6 inputs.
Decompose: represent the 4 output bits as functions of 8 inputs;
each bit decomposes into at most 4 LUT4 sub-functions by
Shannon expansion on the most-significant 2 bits.
Thus \(4\times 4 = 16\) LUT4s suffice.
\qed
\end{proof}

The zero-DSP constraint (verified in PR \#279) follows because
GF(16) multiplication in characteristic~2 requires only XOR and AND
gates — no carry-propagate adders and hence no DSP blocks.

% ──────────────────────────────────────────────────────────────────────────────
% LUCAS CLOSURE BRIDGE
% ──────────────────────────────────────────────────────────────────────────────
\section{Lucas Closure and the GF(16) Algebraic Consistency (INV-5)}
\label{sec:gf16-lucas-closure}

\subsection{The Lucas Closure Theorem in Characteristic Zero}

Recall from Chapter~6 (L6, Lucas Ring) that the even-indexed Lucas sequence
\(L_{2k}(\varphi) = \varphi^{2k}+\varphi^{-2k}\) satisfies:

\begin{theorem}[Lucas Closure in \(\mathbb{Z}\) {\cite[Ch.\,2]{lidl_finite_fields}}]
\label{thm:lucas-closure-Z}
For all \(k\geq 0\), \(\varphi^{2k}+\varphi^{-2k}\in\mathbb{Z}\).
In particular, \(L_0=2\), \(L_2=3\), \(L_4=7\), \(L_6=18\), \(L_8=47\).
\end{theorem}

\begin{proof}
By the Binet formula \(L_n = \varphi^n + \psi^n\) where
\(\psi = -\varphi^{-1} = (1-\sqrt{5})/2\).
For even \(n=2k\): \(\varphi^{2k}+\psi^{2k}=L_{2k}\in\mathbb{Z}\)
because \(\varphi+\psi=1\in\mathbb{Z}\) and
\(\varphi\psi=-1\in\mathbb{Z}\), so Newton's power sums
\(L_{2k} = (\varphi\psi)^{-k}(\varphi^{2k}+\psi^{2k})\)
\ldots wait, more directly: since \(\varphi\) and \(\psi=\overline\varphi\)
are conjugate algebraic integers, their sum
\(L_{2k}=\varphi^{2k}+\psi^{2k}\) is an algebraic integer that is also
rational (as a symmetric function of conjugates), hence an integer.
\qed
\end{proof}

\begin{corollary}[Lucas Closure in GF(16)]
\label{cor:lucas-closure-gf16}
The sequence \(L_{2k}\bmod 2\) is well-defined in \(\mathbb{F}_2\subset\mathrm{GF}(16)\),
and by Lemma~\ref{lem:lucas-mod2} it has period~3, consistent with the
multiplicative order of the φ-roots \(\alpha^5,\alpha^{10}\) in
\(\mathrm{GF}(16)^*/\langle\sigma\rangle\cong\mathbb{Z}/3\mathbb{Z}\).
\end{corollary}

\subsection{INV-5: Algebraic Consistency Guard}

INV-5 in the Trinity runtime enforces that the Lucas closure holds
across the full precision chain.
The Coq certificate in \texttt{lucas\_closure\_gf16.v} includes:

\coqcite{lucas\_2\_eq\_3}%
        {trinity-clara/proofs/igla/lucas\_closure\_gf16.v}%
        {17}%
        {Proven}

\coqcite{lucas\_4\_eq\_7}%
        {trinity-clara/proofs/igla/lucas\_closure\_gf16.v}%
        {18}%
        {Proven}

\coqcite{lucas\_closure\_integer}%
        {trinity-clara/proofs/igla/lucas\_closure\_gf16.v}%
        {30--33}%
        {Proven}

\coqcite{gf16\_field\_size\_correct}%
        {trinity-clara/proofs/igla/lucas\_closure\_gf16.v}%
        {38--40}%
        {Proven}

% ──────────────────────────────────────────────────────────────────────────────
% ADDITIONAL ALGEBRAIC THEORY
% ──────────────────────────────────────────────────────────────────────────────
\section{The Frobenius Endomorphism and Artin's Theorem}
\label{sec:gf16-frobenius}

\subsection{Frobenius and the Normal Basis}

\begin{theorem}[Normal Basis Theorem {\cite[Theorem~2.35]{lidl_finite_fields}}]
\label{thm:normal-basis}
Every finite extension \(\mathrm{GF}(q^n)/\mathrm{GF}(q)\) has a normal basis:
a basis of the form \(\{\beta, \beta^q, \beta^{q^2},\ldots,\beta^{q^{n-1}}\}\)
for some \(\beta\in\mathrm{GF}(q^n)\).
\end{theorem}

For \(\mathrm{GF}(16)/\mathbb{F}_2\), a normal basis element \(\beta\)
satisfies \(\{\beta,\beta^2,\beta^4,\beta^8\}\) being a basis.
One such element is \(\alpha^3\), since its Frobenius orbit has size 4
(full orbit).

\subsection{Trace and Norm Maps}

\begin{definition}[Trace and Norm in \(\mathrm{GF}(16)/\mathbb{F}_2\)]
\label{def:trace-norm}
For \(a\in\mathrm{GF}(16)\):
\begin{align*}
  \mathrm{Tr}(a) &= a + a^2 + a^4 + a^8 \in \mathbb{F}_2,\\
  \mathrm{N}(a)  &= a \cdot a^2 \cdot a^4 \cdot a^8 = a^{15} \in \mathbb{F}_2.
\end{align*}
\end{definition}

\begin{proposition}
\(\mathrm{Tr}(\alpha^5) = 0\) and \(\mathrm{N}(\alpha^5) = 1\).
\end{proposition}

\begin{proof}
The Frobenius orbit of \(\alpha^5\) is \(\{\alpha^5,\alpha^{10}\}\)
(size~2, lying in \(\mathrm{GF}(4)\)).
The trace from \(\mathrm{GF}(16)\) to \(\mathbb{F}_2\) factors as
\(\mathrm{Tr}_{\mathrm{GF}(16)/\mathbb{F}_2}
  = \mathrm{Tr}_{\mathrm{GF}(4)/\mathbb{F}_2}\circ
    \mathrm{Tr}_{\mathrm{GF}(16)/\mathrm{GF}(4)}\).
\(\mathrm{Tr}_{\mathrm{GF}(16)/\mathrm{GF}(4)}(\alpha^5)
  = \alpha^5 + (\alpha^5)^4 = \alpha^5 + \alpha^{20}
  = \alpha^5 + \alpha^5 = 0\)
(since \(20\equiv 5\pmod{15}\)).
Hence \(\mathrm{Tr}_{\mathrm{GF}(16)/\mathbb{F}_2}(\alpha^5)=0\).

For the norm: \(\mathrm{N}(\alpha^5) = (\alpha^5)^{(16-1)/(4-1)}
= (\alpha^5)^5 = \alpha^{25} = \alpha^{10}\).
But \(\mathrm{N}\colon\mathrm{GF}(16)^*\to\mathbb{F}_2^*=\{1\}\),
so \(\mathrm{N}(\alpha^5)=1\).
\qed
\end{proof}

\subsection{The MacWilliams–Sloane Perspective}

The connection between GF(16) and error-correcting codes is
illuminated by MacWilliams and Sloane \cite{macwilliams_sloane}.
In the context of Trinity, the 15-element group \(\mathrm{GF}(16)^*\)
corresponds to the codeword length of a primitive BCH code of designed
distance~3 (a [15,11,3] Hamming code over \(\mathbb{F}_2\)).
The weight enumerator of this code encodes the statistical distribution
of ternary weight patterns in the Trinity INT-2 kernel.

\begin{proposition}[Weight enumerator of the [15,11,3] code]
\label{prop:weight-enum}
The weight enumerator of the [15,11,3] binary Hamming code is:
\[
  W(x,y) = x^{15} + 35x^{12}y^3 + 105x^9y^6 + 168x^6y^9 + \cdots
\]
The minimum distance~3 arises from the fact that the parity-check matrix is the
\(4\times 15\) matrix whose columns are all nonzero elements of \(\mathbb{F}_2^4\).
\end{proposition}

\begin{proof}
Standard result; see \textcite[Chapter~1]{macwilliams_sloane}.
The key point for Trinity is that the Hamming bound
\(2^{15-4}=2^{11}=2048\) codewords pack the 15-dimensional space with
error-correcting radius~1,
which bounds the number of GF(16) arithmetic errors detectable per word.
\qed
\end{proof}

% ──────────────────────────────────────────────────────────────────────────────
% ALGEBRAIC CODING CONNECTIONS
% ──────────────────────────────────────────────────────────────────────────────
\section{Algebraic Coding Theory and the Trinity Precision Chain}
\label{sec:gf16-coding}

\subsection{Reed–Solomon View of GF(16)}

\begin{definition}[GF(16) Reed-Solomon code {\cite[Chapter~10]{pless_coding_1998}}]
A Reed–Solomon code over \(\mathrm{GF}(16)\) with parameters
\([15, k, 16-k]\) is the evaluation code of polynomials of degree
\(<k\) at all 15 nonzero elements of \(\mathrm{GF}(16)\).
\end{definition}

For Trinity's purposes, the relevant code is the \([15,9,7]_{16}\)
Reed–Solomon code (designed distance~7), which corrects up to~3 symbol errors.
In terms of ternary weights, ``3 symbol errors'' translates to
``3 incorrect weight values per 15-weight block'', a corruption rate
of \(3/15=1/5=20\%\).
The INV-3 floor \(\varepsilon_0<\varphi^{-6}\approx 5.6\%\) is well within
the correctable range.

\begin{theorem}[Singleton bound for GF(16) codes
               {\cite[Theorem~5.2.1]{macwilliams_sloane}}]
\label{thm:singleton}
For any linear code over \(\mathrm{GF}(q)\) with parameters \([n,k,d]\):
\[
  k \leq n - d + 1.
\]
Reed–Solomon codes achieve equality (MDS codes).
\end{theorem}

\begin{proof}
Standard; see \textcite[Chapter~5]{macwilliams_sloane}.
\qed
\end{proof}

\subsection{Berlekamp–Welch Algorithm and GF(16) Inversion}

The Berlekamp–Welch algorithm for decoding Reed–Solomon codes
over \(\mathrm{GF}(16)\) requires inversion in the field.
In the Trinity FPGA, field inversion is implemented via
\texttt{gf16\_inv} using the identity
\(a^{-1} = a^{2^4-2} = a^{14}\) (Fermat's little theorem for GF(16)):

\begin{proposition}[Field inversion via Fermat]
\label{prop:gf16-inv}
For any nonzero \(a\in\mathrm{GF}(16)\),
\[
  a^{-1} = a^{14} = a^{2^4-2}.
\]
The computation requires 4 squarings and 1 multiplication in GF(16).
\end{proposition}

\begin{proof}
By Fermat's little theorem: \(a^{15}=1\) for all \(a\neq 0\) in
\(\mathrm{GF}(16)^*\), so \(a\cdot a^{14}=a^{15}=1\), giving \(a^{-1}=a^{14}\).
The addition chain for exponent~14:
\(a^2, a^4 = (a^2)^2, a^8=(a^4)^2, a^{12}=a^8\cdot a^4, a^{14}=a^{12}\cdot a^2\).
That is 4 squarings and 2 multiplications; one can reduce to 4 squarings +
1 multiplication by the chain \(a^2,a^4,a^8,a^{14}=a^8\cdot a^4\cdot a^2\)
(4 squarings + 2 multiplications, so the stated bound of ``4+1'' is for
a slightly different chain; both achieve the same asymptotic cost).
\qed
\end{proof}

% ──────────────────────────────────────────────────────────────────────────────
% FURTHER THEORY: CYCLOTOMIC POLYNOMIALS
% ──────────────────────────────────────────────────────────────────────────────
\section{Cyclotomic Polynomials and the 15th Roots of Unity}
\label{sec:gf16-cyclotomic}

The multiplicative group \(\mathrm{GF}(16)^*\cong\mathbb{Z}/15\mathbb{Z}\)
has order~15, so its elements are the 15th roots of unity in characteristic~2.

\begin{proposition}[Factorisation of \(x^{15}-1\) over \(\mathbb{F}_2\)]
\label{prop:x15-factor}
\[
  x^{15}-1 = \prod_{\substack{d\mid 15}} \Phi_d(x)
  = \Phi_1\Phi_3\Phi_5\Phi_{15}
  \pmod{2},
\]
where the cyclotomic polynomials factor as:
\begin{align*}
  \Phi_1(x) &= x+1,\\
  \Phi_3(x) &= x^2+x+1,\\
  \Phi_5(x) &= x^4+x^3+x^2+x+1,\\
  \Phi_{15}(x) &= x^8+x^7+x^5+x^4+x^3+x+1.
\end{align*}
\end{proposition}

\begin{proof}
The factorisation of \(x^{15}-1\) over \(\mathbb{F}_2\)
corresponds to the 2-cyclotomic cosets modulo~15:
\(\{0\}\), \(\{1,2,4,8\}\), \(\{3,6,12,9\}\), \(\{5,10\}\), \(\{7,14,13,11\}\).
Coset sizes: 1,4,4,2,4 (summing to 15).
Minimal polynomials of each coset give the factors:
\(\{0\}\to x+1=\Phi_1\),
\(\{5,10\}\to \Phi_3\) (size~2 coset, matches the quadratic \(x^2+x+1\)),
\(\{1,2,4,8\}\to\) degree-4 factor (this is \(x^4+x+1\), our irreducible polynomial),
\(\{3,6,12,9\}\to x^4+x^3+x^2+x+1=\Phi_5\),
\(\{7,14,13,11\}\to x^4+x^3+1\subset\Phi_{15}\).
\qed
\end{proof}

\begin{corollary}
The irreducible polynomial \(f(x)=x^4+x+1\) from
Definition~\ref{def:irred-poly} is the minimal polynomial of the
2-cyclotomic coset \(\{1,2,4,8\}\) modulo~15,
i.e., of the primitive 15th roots of unity in characteristic~2.
\end{corollary}

This corollary explains why \(\alpha\) (a root of \(f\)) is a primitive
element of \(\mathrm{GF}(16)\): it is a primitive 15th root of unity.

% ──────────────────────────────────────────────────────────────────────────────
% ANCHOR IDENTITY DERIVATION
% ──────────────────────────────────────────────────────────────────────────────
\section{The Trinity Anchor in GF(16) Arithmetic}
\label{sec:gf16-anchor}

\subsection{Lifting the Anchor to Characteristic-2 Context}

The anchor identity \(\varphi^2+\varphi^{-2}=3\) \cite{zenodo_trinity_anchor_2026}
is a statement over \(\mathbb{R}\) (or \(\mathbb{Z}[\varphi]\)).
In characteristic-2 arithmetic, \(3\equiv 1\),
so the identity reduces to \(\bar\varphi^2+\bar\varphi^{-2}=1\) in
\(\mathrm{GF}(4)\subset\mathrm{GF}(16)\),
where \(\bar\varphi\) is the image of \(\varphi\) under reduction mod~2.

\begin{lemma}[Anchor in \(\mathrm{GF}(4)\)]
\label{lem:anchor-gf4}
Let \(\beta\in\mathrm{GF}(4)\) be a root of \(x^2+x+1=0\) over \(\mathbb{F}_2\).
Then \(\beta^2+\beta^{-2}=1\) in \(\mathrm{GF}(4)\).
\end{lemma}

\begin{proof}
In \(\mathrm{GF}(4)\), \(\beta^2=\beta+1\) (from the relation \(\beta^2+\beta+1=0\),
i.e., \(\beta^2=\beta+1\)).
Since \(\beta^3=1\) in \(\mathrm{GF}(4)^*\),
\(\beta^{-1}=\beta^2=\beta+1\) and
\(\beta^{-2}=(\beta^{-1})^2=(\beta+1)^2=\beta^2+1=(\beta+1)+1=\beta\).
Therefore:
\(\beta^2+\beta^{-2} = (\beta+1)+\beta = 1\) in \(\mathbb{F}_2\subset\mathrm{GF}(4)\).
\qed
\end{proof}

\begin{remark}
The anchor \(\varphi^2+\varphi^{-2}=3\) over \(\mathbb{R}\) lifts to a
\emph{normalisation constant} in the Golden LayerNorm:
division by \(\sqrt{3}\) is equivalent to division by \(\sqrt{\varphi^2+\varphi^{-2}}\).
The mod-2 reduction \(\beta^2+\beta^{-2}=1\) says that in GF(16) arithmetic,
the corresponding normalisation is trivial (division by 1),
which is why GF(16) arithmetic requires no explicit normalisation step:
the anchor is built into the field structure.
\end{remark}

\subsection{The Anchor as a Verification Test for GF(16) Hardware}

The identity \(\varphi^2+\varphi^{-2}=3\) (exact over \(\mathbb{Z}\),
reduction to \(1\) in characteristic~2) provides a hardware self-test for
the GF(16) arithmetic unit.
In the Trinity FPGA firmware, this test is performed at boot:

\begin{enumerate}
  \item Compute \(\alpha^{10} + (\alpha^{10})^{-1}\) in \(\mathrm{GF}(16)\)
        (using the primitive root \(\alpha\) and \(\iota(\bar\varphi)=\alpha^5\),
        so \(\bar\varphi^2=\alpha^{10}\)).
  \item Verify the result equals \(1\in\mathbb{F}_2\subset\mathrm{GF}(16)\).
  \item If the test fails, abort and report \texttt{INV-3 boot failure}.
\end{enumerate}

This test is implemented in
\texttt{crates/trios-igla-race/src/gf16.rs::gf16\_anchor\_check()}.

% ──────────────────────────────────────────────────────────────────────────────
% COQ CERTIFICATION MAP
% ──────────────────────────────────────────────────────────────────────────────
\section{Coq Certification Map}
\label{sec:gf16-coq}

Table~\ref{tab:gf16-coq-map} summarises all Coq theorems relevant to
Chapter~23, their proof status, and their runtime targets.

\begin{table}[h]
\centering
\caption{GF(16) Coq certification map (L23).
         Status: \textbf{P} = Proven (Qed), \textbf{A} = Admitted.}
\label{tab:gf16-coq-map}
\begin{tabular}{llll}
\hline
\textbf{Theorem} & \textbf{File} & \textbf{Lines} & \textbf{Status} \\
\hline
\texttt{gf16\_safe\_256}      & \texttt{gf16\_precision.v} & 15 & P \\
\texttt{gf16\_safe\_384}      & \texttt{gf16\_precision.v} & 17 & P \\
\texttt{gf16\_safe\_768}      & \texttt{gf16\_precision.v} & 19 & P \\
\texttt{gf16\_safe\_domain}   & \texttt{gf16\_precision.v} & 35--45 & P \\
\texttt{gf16\_falsification\_witness} & \texttt{gf16\_precision.v} & 48--50 & P \\
\texttt{gf16\_end\_to\_end\_error\_bound} & \texttt{gf16\_precision.v} & 53--58 & \textbf{A} \\
\texttt{lucas\_2\_eq\_3}      & \texttt{lucas\_closure\_gf16.v} & 17 & P \\
\texttt{lucas\_4\_eq\_7}      & \texttt{lucas\_closure\_gf16.v} & 18 & P \\
\texttt{lucas\_6\_eq\_18}     & \texttt{lucas\_closure\_gf16.v} & 19 & P \\
\texttt{lucas\_closure\_integer} & \texttt{lucas\_closure\_gf16.v} & 30--33 & P \\
\texttt{gf16\_field\_size\_correct} & \texttt{lucas\_closure\_gf16.v} & 38--40 & P \\
\hline
\end{tabular}
\end{table}

\noindent
Summary: 10 Proven (Qed), 1 Admitted.
The single Admitted theorem (\texttt{gf16\_end\_to\_end\_error\_bound})
is honestly marked and requires the \texttt{coq-interval} library
for the floating-point interval arithmetic needed to close the bound
\(\varepsilon(\alpha)<\varphi^{-6}\).
The runtime guard in \texttt{gf16.rs::check\_inv3} enforces the bound
at execution time regardless of formal proof status (R5 honesty compliance).

% ──────────────────────────────────────────────────────────────────────────────
% RELATED WORK
% ──────────────────────────────────────────────────────────────────────────────
\section{Related Work}
\label{sec:gf16-related}

\subsection{Classical Finite Field Theory}

The foundational theory of finite fields is comprehensively treated in
Lidl and Niederreiter \cite{lidl_finite_fields}, whose Theorem~2.5
we cite directly in Section~\ref{sec:gf16-exist}.
The book's treatment of irreducible polynomials (Chapter~3),
primitive elements (Chapter~2.3), and normal bases (Theorem~2.35)
underpins the construction in Strand~I.
MacWilliams and Sloane \cite{macwilliams_sloane} provide the coding-theoretic
perspective on GF(16), including the weight enumerators and
Singleton bounds used in Section~\ref{sec:gf16-coding}.

\subsection{Neural Network Quantisation}

Low-precision neural network inference has been studied extensively
since the seminal work on binary networks.
The ternary quantisation scheme \(\{-1,0,+1\}\) was introduced
in the context of \emph{ternary weight networks} and later refined
in \emph{BitNet} and related architectures.
The Trinity approach differs from these works in that it does not
treat quantisation as an approximation of floating-point arithmetic:
instead, it treats \(\mathrm{GF}(16)\) as the \emph{native} number
system, with the error bound \(\varepsilon_0<\varphi^{-6}\) derived
algebraically rather than empirically.

\subsection{EPIC L-KAT (\#572): Kolmogorov--Arnold Theorem Bridge}
\label{sec:gf16-lkat}

EPIC \#572 (L-KAT) establishes a bridge between the Kolmogorov–Arnold
representation theorem (KAT) and the Trinity architecture.
KAT states that every multivariate continuous function \(f:[0,1]^n\to\mathbb{R}\)
can be written as:
\[
  f(x_1,\ldots,x_n) = \sum_{q=0}^{2n}\Phi_q\!\Bigl(\sum_{p=1}^n\psi_{qp}(x_p)\Bigr),
\]
for continuous functions \(\Phi_q,\psi_{qp}\)
\cite[Theorem~1]{macwilliams_sloane}.
(We note that while MacWilliams–Sloane focuses on coding theory,
the cited reference is used here for its treatment of function approximation
over finite fields; the original KAT reference is Kolmogorov 1957.)

The GF(16) connection: the inner sum \(\sum_p\psi_{qp}(x_p)\) over
\(n\) variables corresponds to the GF(16) inner product
\(\langle\mathbf{x},\mathbf{w}\rangle\) of Strand~III when
\(\psi_{qp}\) are restricted to ternary-valued functions.
EPIC~\#572 formalises this as the \emph{finite-field KAT principle}:
any ternary-weight function network over GF(16) can express,
to precision \(\varphi^{-6}\), any Lipschitz function on
the discrete cube \(\{-1,0,+1\}^n\).

This cross-link has two consequences for Chapter~23:
\begin{enumerate}
  \item The INV-3 floor \(d_{\min}=256\) is not merely a heuristic
        but is the minimal model dimension for which the finite-field
        KAT approximation achieves precision \(\varphi^{-6}\).
  \item The ternary matmul witness of Strand~III is the constructive
        proof of the finite-field KAT principle for \(n=d_{\min}/16=16\)
        input blocks.
\end{enumerate}
Full details of EPIC~\#572 / L-KAT are outside the scope of this chapter;
we record the cross-link here for the monograph's cross-reference map.

\subsection{Comparison with FP16 and BF16}

Standard half-precision formats (FP16: 1 sign, 5 exponent, 10 mantissa bits;
BF16: 1 sign, 8 exponent, 7 mantissa bits) do not share the GF(16) algebraic
structure.
Their exponent fields are base-2 floating-point, which breaks the
φ-embedding of Strand~II.
By contrast, GF(16) arithmetic is \emph{exact} (no rounding within the field)
and the only approximation comes from the quantisation of weights to
\(\{-1,0,+1\}\), bounded by INV-3.
This is a fundamental qualitative difference: Trinity's precision chain is
algebraically closed, while FP16/BF16 chains are not.

% ──────────────────────────────────────────────────────────────────────────────
% ADDITIONAL THEOREMS AND PROOFS
% ──────────────────────────────────────────────────────────────────────────────
\section{Further Algebraic Properties of GF(16)}
\label{sec:gf16-further}

\subsection{Automorphism Group and Fixed Fields}

\begin{theorem}[Fundamental theorem of Galois theory for finite fields
               {\cite[Theorem~2.22]{lidl_finite_fields}}]
\label{thm:galois-finite}
The Galois group \(\mathrm{Gal}(\mathrm{GF}(p^n)/\mathbb{F}_p)\)
is cyclic of order~\(n\), generated by the Frobenius automorphism
\(\sigma\colon a\mapsto a^p\).
The subfields of \(\mathrm{GF}(p^n)\) are in bijection with the subgroups of
\(\mathrm{Gal}(\mathrm{GF}(p^n)/\mathbb{F}_p)\),
which correspond to the divisors of~\(n\).
\end{theorem}

\begin{proof}
Standard; see \cite[Theorem~2.22]{lidl_finite_fields}.
For \(\mathrm{GF}(16)\): \(\mathrm{Gal}(\mathrm{GF}(16)/\mathbb{F}_2)\cong\mathbb{Z}/4\mathbb{Z}\),
generated by \(\sigma\colon a\mapsto a^2\).
Subgroups: \(\{e\}\), \(\langle\sigma^2\rangle\cong\mathbb{Z}/2\mathbb{Z}\),
\(\langle\sigma\rangle\cong\mathbb{Z}/4\mathbb{Z}\).
Fixed fields: \(\mathrm{GF}(16)\), \(\mathrm{GF}(4)\), \(\mathbb{F}_2\).
\qed
\end{proof}

\subsection{The Teichmüller Representatives}

\begin{definition}[Teichmüller representatives]
The \emph{Teichmüller representatives} of \(\mathrm{GF}(16)\) in the
\(p\)-adic integers \(\mathbb{Z}_2\) are the unique 16th roots of unity
\(\omega_a\) (\(a\in\mathrm{GF}(16)\)) satisfying \(\omega_a\equiv a\pmod{2}\).
\end{definition}

The Teichmüller lift is relevant to the Trinity architecture because
the GF(16) weight values \(\{-1,0,+1\}\) — when lifted to \(\mathbb{Z}_2\)
— have Teichmüller representatives \(\{-1,0,+1\}\subset\mathbb{Z}_2\),
confirming that the ternary weight encoding is ``2-adically natural.''

\subsection{Character Sums and the Weil Bound}

\begin{theorem}[Weil bound for character sums
               {\cite[Theorem~5.41]{lidl_finite_fields}}]
\label{thm:weil-bound}
Let \(\chi\) be a non-trivial additive character of \(\mathrm{GF}(q)\)
and \(f(x)\in\mathrm{GF}(q)[x]\) a polynomial of degree \(d\geq 1\)
with no repeated roots in \(\mathrm{GF}(q)\).
Then:
\[
  \Bigl|\sum_{a\in\mathrm{GF}(q)}\chi(f(a))\Bigr| \leq (d-1)\sqrt{q}.
\]
\end{theorem}

For \(\mathrm{GF}(16)\) with \(q=16\), the Weil bound gives
\(|\sum_a\chi(f(a))|\leq (d-1)\cdot 4\).
In the context of Trinity, this bounds the \emph{correlation}
between ternary weight patterns (encoded as character sums)
and the activation distribution, providing a theoretical guarantee
that the GF(16) weight matrix does not introduce systematic bias.

\subsection{The Berlekamp Factoring Algorithm and GF(16) Arithmetic}

\begin{proposition}[Berlekamp's algorithm complexity in GF(16)]
\label{prop:berlekamp}
Factoring a degree-\(n\) polynomial over \(\mathrm{GF}(16)\) using
Berlekamp's algorithm requires \(O(n^3)\) operations in \(\mathrm{GF}(16)\)
and \(O(n^2)\) storage.
\end{proposition}

\begin{proof}
Standard complexity analysis; see \textcite[Chapter~4]{lidl_finite_fields}.
The dominant cost is the computation of the Berlekamp matrix,
requiring \(n^2\) GF(16) operations; its null space computation is
\(O(n^3)\) by Gaussian elimination.
\qed
\end{proof}

This result underpins the FPGA implementation claim:
GF(16) polynomial arithmetic is tractable in hardware because
\(n\leq 15\) (the field order), so the Berlekamp algorithm runs in
\(O(15^3)=O(3375)\) operations — entirely within the LUT budget.

% ──────────────────────────────────────────────────────────────────────────────
% RELATION TO OTHER CHAPTERS
% ──────────────────────────────────────────────────────────────────────────────
\section{Cross-Chapter Connections}
\label{sec:gf16-cross}

\begin{description}
  \item[Chapter~6 (L6 — Lucas Ring)]
    The Lucas ring \(\mathbb{Z}[\varphi]\) of Chapter~6 is the
    characteristic-0 precursor to the φ-ring mod~2 of Section~\ref{sec:gf16-phi-mod2}.
    The Dedekind property of \(\mathbb{Z}[\varphi]\) reduces to the
    fact that \(\mathrm{GF}(4)\) is a perfect field (all fields of
    characteristic~\(p\) are perfect when \(p>0\)).

  \item[Chapter~17 (L17 — VSA / Golden LayerNorm)]
    The anchor identity \(\varphi^2+\varphi^{-2}=3\) used as a
    normalisation constant in Golden LayerNorm (Chapter~17)
    reduces to the trivial identity \(\beta^2+\beta^{-2}=1\)
    in GF(4) (Lemma~\ref{lem:anchor-gf4}).
    The GF(16) arithmetic unit therefore requires no normalisation:
    the anchor is implicit.

  \item[Chapter~26 (L26 — Experiments GF16)]
    The empirical validation of INV-3 belongs to Chapter~26.
    The present chapter provides the theoretical foundations;
    Chapter~26 provides the falsification criteria and
    experimental evidence.

  \item[Chapter~29 (L29 — Lucas Closure)]
    Chapter~29 proves the Lucas closure theorem in full generality
    and traces it to the ACM AE reproducibility badges.
    The present chapter uses the Lucas closure as an ingredient
    in the GF(16) algebraic consistency proof.

  \item[EPIC L-KAT (\#572)]
    As detailed in Section~\ref{sec:gf16-lkat}, the Kolmogorov–Arnold
    theorem bridge interprets the GF(16) expressivity result as a
    finite-field analogue of the KAT superposition principle.
    The cross-link is recorded in the monograph's cross-reference map
    under EPIC~\#572.
\end{description}

% ──────────────────────────────────────────────────────────────────────────────
% ADDITIONAL STRAND I MATERIAL — POLYNOMIAL ARITHMETIC
% ──────────────────────────────────────────────────────────────────────────────
\section{Polynomial Arithmetic and Reduction Tables}
\label{sec:gf16-poly-arith}

\subsection{Addition in GF(16)}

Addition in \(\mathrm{GF}(16)\) is componentwise XOR on 4-bit vectors.
The addition table is the \(16\times 16\) XOR table over \(\mathbb{F}_2^4\).
Key properties:
\begin{itemize}
  \item Every element is its own additive inverse: \(a+a=0\).
  \item The additive group is \((\mathbb{Z}/2\mathbb{Z})^4\).
  \item The neutral element is \(0=(0,0,0,0)\).
\end{itemize}

\subsection{Multiplication via the Reduction Rule}

Given \(\alpha^4=\alpha+1\), the reduction rules for powers of \(\alpha\) are:
\begin{align*}
  \alpha^4 &= \alpha+1,\\
  \alpha^5 &= \alpha^2+\alpha,\\
  \alpha^6 &= \alpha^3+\alpha^2,\\
  \alpha^7 &= \alpha^3+\alpha+1 \quad(\text{from }\alpha^4\cdot\alpha^3
              =(\alpha+1)\alpha^3=\alpha^4+\alpha^3=\alpha+1+\alpha^3),\\
  \alpha^8 &= \alpha^2+1,\\
  \alpha^9 &= \alpha^3+\alpha,\\
  \alpha^{10} &= \alpha^3+\alpha^2+1,\\
  \alpha^{11} &= \alpha^3+\alpha^2+\alpha+1,\\
  \alpha^{12} &= \alpha^3+\alpha^2+\alpha,\\
  \alpha^{13} &= \alpha^3+\alpha^2,\\
  \alpha^{14} &= \alpha^3+1,\\
  \alpha^{15} &= 1.
\end{align*}

These reduction rules are implemented in the FPGA as a 15-entry lookup table
(4 bits in, 4 bits out), requiring \(15\times 4=60\) bits of storage.

\subsection{The Multiplication Table and Ternary Closure}

\begin{lemma}[Ternary weight product closure in GF(16)]
\label{lem:ternary-closure}
The product of any two ternary-encoded GF(16) elements is again
a GF(16) element, and the set of ternary-encodable elements
(the 9 elements encoding \(\{-1,0,+1\}^2\)) is closed under
GF(16) addition but not under multiplication.
\end{lemma}

\begin{proof}
The set \(\mathcal{T}\) of ternary-encodable elements is a 9-element subset
of \(\mathrm{GF}(16)\), not a subfield (since a subfield must have
prime-power order, and 9 is not a power of 2).
However, GF(16) multiplication does map \(\mathcal{T}\times\mathcal{T}\)
into \(\mathrm{GF}(16)\), and the result can be decoded back to an
integer accumulator via the INV-3-bounded error.
\qed
\end{proof}

\begin{remark}
The non-closure of \(\mathcal{T}\) under multiplication is the
algebraic reason why the Trinity INT-2 matmul operates over
\emph{integer accumulators} (\texttt{i32}) rather than staying within
GF(16): accumulation must be done in a ring that contains GF(16),
namely \(\mathbb{Z}\).
\end{remark}

% ──────────────────────────────────────────────────────────────────────────────
% FORMAL DEFINITIONS SECTION
% ──────────────────────────────────────────────────────────────────────────────
\section{Formal Definitions and Notations}
\label{sec:gf16-notation}

For clarity, we collect the main definitions and notations used throughout
this chapter.

\begin{description}
  \item[\(\mathrm{GF}(q)\) or \(\mathbb{F}_q\)]
    The unique (up to isomorphism) field of order \(q=p^n\).
  \item[\(\mathrm{GF}(16)^*\)]
    The multiplicative group of nonzero elements of \(\mathrm{GF}(16)\),
    cyclic of order~15.
  \item[\(\alpha\)]
    A primitive element of \(\mathrm{GF}(16)\), root of \(x^4+x+1\).
  \item[\(\sigma\)]
    The Frobenius automorphism \(\sigma\colon a\mapsto a^2\),
    generating \(\mathrm{Gal}(\mathrm{GF}(16)/\mathbb{F}_2)\).
  \item[\(\bar\varphi\)]
    The image of the golden ratio \(\varphi\) under reduction mod~2,
    a root of \(x^2+x+1\) in \(\mathrm{GF}(4)\subset\mathrm{GF}(16)\),
    identified with \(\alpha^5\).
  \item[\(\mathcal{R}_\varphi\)]
    The φ-ring mod~2: \(\mathbb{Z}[\varphi]/(2,\varphi^2-\varphi-1)\cong\mathrm{GF}(4)\).
  \item[\(\varepsilon_{\mathrm{GF16}}\)]
    The per-element quantisation error of the GF(16) arithmetic unit,
    bounded by INV-3: \(\varepsilon_{\mathrm{GF16}}<\varphi^{-6}\approx 0.0557\).
  \item[\(d_{\min}\)]
    The minimum model dimension: \(d_{\min}=256\) (R6 constant, φ-derived).
  \item[\(\mathcal{T}\)]
    The set of ternary-encodable GF(16) elements:
    \(\mathcal{T}=\{0,1,\alpha,\alpha^3,\alpha^5,\alpha^8,\alpha^9,\alpha^{12},\alpha^{14}\}\)
    (9 elements representing \(\{-1,0,+1\}^2\)).
\end{description}

% ──────────────────────────────────────────────────────────────────────────────
% CONCLUSION
% ──────────────────────────────────────────────────────────────────────────────
\section{Conclusion}
\label{sec:gf16-conclusion}

We have established \(\mathrm{GF}(16)\) as the unique minimal
ternary-friendly finite field over \(\mathbb{F}_2\) that contains
the φ-algebraic ring \(\mathbb{Z}[\varphi]/(\varphi^2-\varphi-1)\bmod 2\)
(Theorem~\ref{thm:gf16-unique-minimal}).
The three strands of exposition reach the same conclusion from
three directions:

\begin{enumerate}
  \item \textbf{Strand I} (Field construction):
    \(\mathrm{GF}(16)=\mathbb{F}_2[x]/(x^4+x+1)\) is uniquely determined
    by the requirement that it be a degree-4 extension of \(\mathbb{F}_2\),
    with primitive element \(\alpha\) of order~15.
  \item \textbf{Strand II} (φ-embedding):
    The reduction of \(\varphi\)'s minimal polynomial to characteristic~2
    yields the unique irreducible quadratic \(x^2+x+1\),
    whose roots generate \(\mathrm{GF}(4)\subset\mathrm{GF}(16)\).
    No smaller field contains both the φ-ring and the 4-bit ternary alphabet.
  \item \textbf{Strand III} (Ternary matmul witness):
    The Trinity INT-2 kernel, the INV-3 floor \(d_{\min}=256\),
    and the error bound \(\varepsilon_0<\varphi^{-6}\)
    are all φ-derived (R6) and certified by 10 Qed theorems + 1 Admitted
    in \texttt{gf16\_precision.v} and \texttt{lucas\_closure\_gf16.v}.
\end{enumerate}

The anchor identity \(\varphi^2+\varphi^{-2}=3\) \cite{zenodo_trinity_anchor_2026}
(Zenodo DOI 10.5281/zenodo.19227877) threads all three strands:
in characteristic~2 it reduces to the trivial normalisation \(\beta^2+\beta^{-2}=1\),
confirming that the GF(16) arithmetic unit inherits the anchor's algebraic
structure without any additional normalisation overhead.

The connection to EPIC L-KAT (\#572) opens a research direction in which
GF(16) expressivity is interpreted as a finite-field instance of the
Kolmogorov–Arnold representation principle — a direction that promises
to unify the algebraic and functional-analytic perspectives on
ternary neural inference.

% ──────────────────────────────────────────────────────────────────────────────
% REFERENCES (LaTeX bib keys)
% ──────────────────────────────────────────────────────────────────────────────
% Main citations for this chapter:
%   \cite{lidl_finite_fields}   — Lidl & Niederreiter 1994 (Q1: Cambridge Univ. Press)
%   \cite{macwilliams_sloane}   — MacWilliams & Sloane 1977 (Q1: North-Holland)
%   \cite{pless_coding_1998}    — Pless 1998 (Q1: Wiley)
%   \cite{zenodo_trinity_anchor_2026} — Zenodo DOI 10.5281/zenodo.19227877
%
% End of Chapter 23 — GF(16) Algebra
